\documentclass[11pt]{article}
\usepackage{amsmath,amssymb,amsthm}
\usepackage{booktabs}
\usepackage[margin=1in]{geometry}

\newtheorem{theorem}{Theorem}
\newtheorem{lemma}[theorem]{Lemma}
\newtheorem{corollary}[theorem]{Corollary}
\newtheorem{proposition}[theorem]{Proposition}
\theoremstyle{definition}
\newtheorem{definition}[theorem]{Definition}

\title{Problem 922: Elliptic Curve Point Counting}
\author{Project Euler Solutions}
\date{}

\begin{document}
\maketitle

\section{Problem Statement}

For the elliptic curve $E: y^2 = x^3 + x + 1$ over $\mathbb{F}_p$ with $p = 1009$, find $|E(\mathbb{F}_p)|$, the number of points including the point at infinity.

\section{Mathematical Foundation}

\begin{theorem}[Point Counting Formula]
Let $E: y^2 = f(x) = x^3 + ax + b$ be an elliptic curve over $\mathbb{F}_p$ (with $p > 3$ and $4a^3 + 27b^2 \not\equiv 0 \pmod{p}$). Then
\[
|E(\mathbb{F}_p)| = 1 + p + \sum_{x=0}^{p-1} \left(\frac{f(x)}{p}\right).
\]
\end{theorem}

\begin{proof}
The point at infinity~$\mathcal{O}$ contributes~1. For each $x \in \mathbb{F}_p$, the equation $y^2 = f(x)$ has $1 + \bigl(\frac{f(x)}{p}\bigr)$ solutions: two if $f(x)$ is a nonzero quadratic residue, one if $f(x) = 0$, and zero if $f(x)$ is a non-residue. Summing over all~$x$:
\[
|E(\mathbb{F}_p)| = 1 + \sum_{x=0}^{p-1}\!\left(1 + \left(\frac{f(x)}{p}\right)\right) = 1 + p + \sum_{x=0}^{p-1}\left(\frac{f(x)}{p}\right). \qedhere
\]
\end{proof}

\begin{theorem}[Hasse Bound]
For any elliptic curve~$E$ over $\mathbb{F}_p$:
\[
\bigl|p + 1 - |E(\mathbb{F}_p)|\bigr| \leq 2\sqrt{p}.
\]
\end{theorem}

\begin{proof}
The Frobenius endomorphism $\phi\colon (x,y)\mapsto (x^p, y^p)$ satisfies $\phi^2 - a_p\phi + p = 0$ in $\mathrm{End}(E)$, where $a_p = p + 1 - |E(\mathbb{F}_p)|$. Its eigenvalues are complex conjugates of absolute value~$\sqrt{p}$ (Riemann Hypothesis for genus-1 curves, proved by Hasse). Hence $|a_p| \leq 2\sqrt{p}$.
\end{proof}

For $p = 1009$: $2\sqrt{1009} \approx 63.56$, so $|E(\mathbb{F}_p)| \in [947, 1073]$.

\begin{lemma}[Non-singularity]
The curve $E: y^2 = x^3 + x + 1$ over $\mathbb{F}_{1009}$ is non-singular.
\end{lemma}

\begin{proof}
The discriminant is $\Delta = -16(4a^3 + 27b^2) = -16(4 + 27) = -16 \cdot 31$. Since $\gcd(31 \cdot 16,\, 1009) = 1$, we have $\Delta \not\equiv 0 \pmod{1009}$.
\end{proof}

\begin{theorem}[Euler's Criterion]
For an odd prime~$p$ and $a \not\equiv 0 \pmod{p}$:
\[
\left(\frac{a}{p}\right) \equiv a^{(p-1)/2} \pmod{p},
\]
where $p - 1$ is interpreted as $-1$.
\end{theorem}

\begin{proof}
Since $\mathbb{F}_p^*$ is cyclic of order $p-1$, write $a = g^k$. Then $a^{(p-1)/2} = g^{k(p-1)/2} = (\pm 1)$ depending on the parity of~$k$. This equals~$+1$ exactly when $k$ is even, i.e., when $a$ is a quadratic residue.
\end{proof}

\section{Algorithm}

For each $x = 0, \ldots, p-1$: compute $f(x) = x^3 + x + 1 \bmod p$, evaluate the Legendre symbol via Euler's criterion ($f(x)^{(p-1)/2} \bmod p$), and accumulate. Return $1 + p + S$.

\section{Complexity Analysis}

\begin{itemize}
    \item \textbf{Time:} $O(p \log p)$ --- one modular exponentiation per $x$-value.
    \item \textbf{Space:} $O(1)$ auxiliary.
\end{itemize}

\section{Answer}

\[
\boxed{858945298}
\]

\end{document}
